% The same 4x4 maze twice: as a drawing of walls, and as the graph this report
% actually works on. Both panels are generated from the one edge set listed in
% Section 2, so they cannot drift apart.
\begin{figure}[htbp]
  \centering
  \begin{subfigure}[b]{0.46\textwidth}
    \centering
    \begin{tikzpicture}[scale=0.95]
      \fill[cellgrey] (0,0) rectangle (4,-4);
      % outer boundary, with the exit opening left out of the left side
      \draw[wall segment] (0,0) -- (4,0) -- (4,-4) -- (0,-4);
      \draw[wall segment] (0,-1) -- (0,-4);
      % walls between cells that no door connects
      \draw[wall segment] (2,0)  -- (2,-1);
      \draw[wall segment] (1,-1) -- (1,-2);
      \draw[wall segment] (3,-1) -- (3,-2);
      \draw[wall segment] (2,-2) -- (2,-3);
      \draw[wall segment] (2,-3) -- (2,-4);
      \draw[wall segment] (0,-3) -- (1,-3);
      \draw[wall segment] (1,-2) -- (2,-2);
      \draw[wall segment] (2,-3) -- (3,-3);
      \draw[wall segment] (3,-1) -- (4,-1);
      % the single exit
      \draw[doorgreen, -{Stealth[length=3mm]}, line width=1.2pt]
            (0.5,-0.5) -- (-0.8,-0.5);
      \node[exitorange, font=\scriptsize\bfseries, left] at (-0.85,-0.5) {exit};
    \end{tikzpicture}
    \caption{What Minos is shown: walls, one way out.}
    \label{fig:maze-as-walls}
  \end{subfigure}\hfill
  \begin{subfigure}[b]{0.46\textwidth}
    \centering
    \begin{tikzpicture}[scale=1.27]
      \foreach \x in {0,1,2,3} \foreach \y in {0,1,2,3}
        \node[node vertex] (v\x\y) at (\x,-\y) {};
      \node[exit vertex] at (0,0) {};   % overdrawn: the one exit
      \foreach \a/\b in {v00/v10, v20/v30, v11/v21, v02/v12, v22/v32,
                         v03/v13, v23/v33,
                         v00/v01, v01/v02, v10/v11, v12/v13,
                         v20/v21, v21/v22, v31/v32, v32/v33}
        \draw[door] (\a) -- (\b);
      \node[exitorange, font=\scriptsize\bfseries, left] at (-0.35,0) {exit};
    \end{tikzpicture}
    \caption{What the algorithm works on: a spanning tree of the grid graph.}
    \label{fig:maze-as-graph}
  \end{subfigure}
  \caption{A maze is not a picture but a graph. Cells become vertices, doors
  become edges, and a wall is simply an edge of the grid graph that is
  absent. The two panels show the same $4\times4$ maze.}
  \label{fig:maze-and-its-graph}
\end{figure}
